\section{Introduction}
\label{sec:introduction}
Text-to-motion generation, which synthesizes diverse and temporally coherent human motions from natural language descriptions, poses a significant challenge in computer vision, graphics, and robotics. Despite notable advancements~\citep{rombach2022high,gpt3,gpt4, yang2025sigman, li2024dispose, chen2025dancetogether, sora, cheng2024learning, xu2023survey} in large-scale generative models for text, images, 3D, and videos—showcasing exceptional zero-shot capabilities—the motion domain~\citep{mdm,t2mgpt,zhou2023let,actformer,momask} remains considerably behind. This discrepancy is attributed not to a lack of algorithmic innovation, but to inherent limitations in data scale and model architecture that hinder robust generalization, thereby constraining real-world applicability.

To address the generalization challenges faced by state-of-the-art methods such as MDM~\citep{mdm}, MotionGPT~\citep{motiongpt}, and MoMask~\citep{momask}, which are limited by dataset constraints in HumanML3D~\citep{humanml3d} and MotionX~\citep{motionx}, earlier approaches~\citep{motionclip, avatarclip} aligned motion sequences with the image embedding space of CLIP~\citep{clip} using rendered frames. While these methods somewhat alleviate data scarcity, they suffer from inherent modality mismatches, as static image embeddings fail to capture temporal dynamics, resulting in incoherent motions and limited compositional reasoning. Recent efforts~\citep{lu2024scamo, xu2024motionbank, omg, wang2024quo} have aimed to enhance generalization by scaling up motion generation models with larger datasets (e.g., 250 hours of motion in ScaMo~\cite{lu2024scamo} and one billion parameters in OMG~\cite{omg}). These methods exhibit improved motion diversity and novel language alignment compared to previous methods~\cite{mdm, motiongpt, momask}, which were trained on more restricted datasets~\cite{humanml3d, motionx}. Nonetheless, impeded by constraints related to model capacity and the inherent limitations of the datasets (\eg quality, diversity, and size), the full potential of this scaling-up formulation remains largely unexploited, particularly when these text inputs entail long-term motion compositions and complex descriptions.

Therefore, we contend that achieving human-level motion generation necessitates a paradigm shift akin to the \emph{``scaling hypothesis"}: sufficiently large and diverse \fk{high quality} training data, combined with scaled model architectures, can unlock emergent zero-shot capabilities\fk{, especially performing the complex compositional motions}. 
\fk{To this end, in this paper we explore three key components for zero-shot motion generation: 1) large-scale and high-quality motion dataset, 2) scalable and model architecture, 3) effective evaluation benchmark.}

\fk{Firstly,} we introduce MotionMillion, a large-scale motion dataset with comprehensive text annotations. We propose a novel and efficient 
\fk{motion annotation mechanism, including motion descriptions and high quality motion capturing,} sourced from web-scale human motion videos. Our framework autonomously harvests human motion from unlabeled videos through kinematic regression, generates semantically rich captions using advanced vision-language models (\eg, gpt-4o~\cite{gpt4}), and implements a multi-stage filtering process to eliminate scene cuts and static pose jitters. This meticulous curation results in a dataset comprising \fk{over} \textbf{2000 hours} of high-quality text-motion pairs, encompassing over \textbf{2 million} motion sequences—\textbf{20 times} larger than existing resources. By unifying annotations across extant datasets (\eg, HumanML3D~\citep{humanml3d}, MotionX~\cite{motionx}), we establish a temporally coherent and compositionally diverse foundation for scaling. 

Leveraging this dataset, we \fk{further explore the effective scalable model to accommodate our large-scale dataset. LLAMA, as a transformer decoder-only model, has significantly demonstrated its scalability in text generation tasks. In the motion field, it has been proven in Scamo~\citep{lu2024scamo} that a scaling law curve can be drawn with the increase in data volume. Therefore, 1) we first use Finite Scalar Quantization (FSQ) as \textbf{Efficient Motion Tokenization:} to discretely encode the motion data. This is a more stable and efficient way compared to VQ-VAE. However, we found that although the discretization method of FSQ can effectively encode motion data when the data scale is limited, due to the extremely large scale of our data, directly using the FSQ model will cause jitter in the reconstructed motion. We believe that this is because the information loss caused by the discretization method of FSQ becomes more serious as the data scale increases, leading the model to wrongly model high-frequency information. To solve this problem, we propose to use wavelet transformation to preprocess the motion data before inputting it into the reconstruction network. After obtaining the decoder outputs, we further utilize the inverse wavelet transformation to obtain the final reconstructed motions. By this means, we can effectively encode the motion data and reduce the jitter phenomenon caused by discretive information loss. After completing the discrete compression encoding, we further utilize the LLAMA architecture to implement}
2) \textbf{Scalable Motion Generation:} We
design a bidirectional transformer that jointly models text-motion cross-attention and autoregressive motion token prediction, enabling compositional motion synthesis. Starting from a 1B parameter base, we progressively scale model depth to the final 7B scale, observing emergent zero-shot capabilities. As shown in Fig.~\ref{fig:teaser}, our model could deal with 
\fk{various texts, especially can follow complex long texts.}

To systematically assess the zero-shot generalization capabilities of models, we \fk{further} introduce MotionMillion-Eval, a new benchmark comprising 126 diverse prompts across 7 categories, ranging from daily life scenarios to inhuman motions. Our evaluation focuses on three key aspects: text-motion alignment, motion smoothness, and physical feasibility of motions. Our findings indicate that the 7 billion parameter model, MotionMillion, successfully trained on the MotionMillion dataset, exhibits robust zero-shot generalization abilities. This advancement paves the way for advancing the motion generation task towards zero-shot applications. 

Our contributions can be summarized as follows:
\begin{itemize}
    \item We propose a high-quality annotation pipeline of human motions from video data and build MotionMillion, a large-scale human motion dataset, which is currently the largest human motion with the highest quality, and its scale and diversity drive the research in the field of human motion towards zero-shot application.
    \item We \fk{propose to leverage the wavelet transformation to descrease the jitter phenomonen from FSQ, and we further scale our model to} 7B parameters via an effective scalable architecture, demonstrating strong generalization for out-of-domain complex compositional motions.
    \item We built the MotionMillion-Eval benchmark according to industry standards, which is the first proposed evaluation that can be used for zero-shot capability verification.
\end{itemize}



